\documentclass[11pt,twoside,letter]{article}
\usepackage[utf8]{inputenc}
\usepackage{amsmath}
\usepackage{amsfonts}
\usepackage{amssymb}
\usepackage[spanish,es-noshorthands]{babel}
\usepackage[T1]{fontenc}
\usepackage{lmodern}
\usepackage{graphicx,hyperref}
\usepackage{tikz,pgf}
\usepackage{multicol}
\usepackage{subfig}
\usepackage[width=7.5in,height=9.5in]{geometry}
\usepackage{fancyhdr}
\pagestyle{fancy}
\fancyhead[LE]{\includegraphics[height=12pt]{Images/logo-colegio.png} Matemáticas $11^{\circ}$}
\fancyhead[RE]{}
\fancyhead[RO]{\textit{Germ\'an Avenda\~no Ram\'irez, Lic. U.D., M.Sc. U.N.}}
\fancyhead[LO]{}

\author{Germ\'an Avenda\~no Ram\'irez, Lic. U.D., M.Sc. U.N.}
\title{\begin{minipage}{.2\textwidth}
\includegraphics[height=1.75cm]{Images/logo-colegio.png}\end{minipage}
\begin{minipage}{.55\textwidth}
\begin{center}
Animaplano 03, $9^{\circ}$  \\
Matemáticas $11^{\circ}$
\end{center}
\end{minipage}\hfill
\begin{minipage}{.2\textwidth}
\includegraphics[height=1.75cm]{Images/logo-sed.png} 
\end{minipage}}
\date{}
\begin{document}
\maketitle
Nombre: \hrulefill Curso: \underline{\hspace*{44pt}} Fecha: \underline{\hspace*{2.5cm}}
\section*{Cuestionario}
\begin{enumerate}
 \item Reste $4^{2}$, del triple de 30
 \item $\frac{5}{8}$, de un ángulo de $120^{\circ}$
 \item En años, $\frac{1}{2}$ década + 12 lustros:
 \item El triple de la mitad de 44
 \item En $y=75x+3$, la pendiente es: ("recuerde que $m$ es la pendiente de una recta cuya ecuaci\'{o}n es $y=mx+b$")
 \item Si $y-x=25$ para $x=51$, luego $y=$?
 \item El doble de $7^{2}$
 \item El área de un rectángulo de perímetro 40 y cuyo largo es 6 unidades más que el ancho.
 \item $(2^{2})^{2}\cdot (2^{4}\div 2^{2})+2^{3}+2=$
 \item Perímetro de un pentágono regular de lado 17 cm
 \item Décimo noveno número primo
 \item Convierta 5700 dm (decímetros) a Dm (Decámetros). Un decámetro (Dm) equivale a 10m y cada metro equivale a 10 decímetros (dm).
 \item $\sqrt{400}+\sqrt{121}+\sqrt{25}=$
 \item Si $a=b$, $a\cdot b=169$, luego $2b=$?
 \item Los meses en $\frac{2}{5}$ de un lustro:
 \item $\frac{1}{3}$ de $\frac{1}{3}$ de 297
 \item $\sqrt{64}\cdot \sqrt{8^{2}}=$
 \item En $y=74x-13$, la pendiente es:
 \item Dos rectas son paralelas si $m_{1}=m_{2}$ V:73 \quad F:85
 \item 1/3 del triple de 61
 \item El complemento de $49^{\circ}$. (Dos ángulos son complementarios si suman $90^{\circ}$)
 \item La media aritmética (promedio) entre 24, 20, 22, 23, 21
 \item El sexto número primo
 \item $\log_{2}256+\log_{3}81+\log100=$ (Recuerde que $\log_{a}x=n$ s\'{i} y s\'{o}lo s\'{i} $a^{n}=x$)
 \item $\sqrt[3]{216}=$
 \item En horas, 1/3 de dos días:
 \item Pendiente de la recta que pasa por los puntos $\left(\frac{3}{2},2\right)$ y $(3,8)$. (La pendiente de la recta que pasa por los puntos $(x_{1},y_{1})$ y $(x_{2},y_{2})$ está dada por $m=\dfrac{y_{2}-y_{1}}{x_{2}-x_{1}}$)
 \item $\{(-28)\div(-4)\}\cdot(-2)+28$
 \item La distancia entre los puntos A(2,3) y B(10,18) es?
 \item El valor de la abscica en el punto (28,18)
 \item Los 3/5 de 1/3 de 135
 \item Si $y=\frac{x}{3}+2$ y $x=105$, entonces $y=$?
 \item 1/2 de un número, sumado con 1/3 de 99, da 62. ¿El número es?
 \item El décimo término de la progresión aritmética --4, 4, 12, \ldots
 \item $4^{3}+\log_{5}125=$
\end{enumerate}
\section*{Animaplano}
\begin{center}
\begin{tikzpicture}
 \fill (0,0) node[above]{0} circle (0.2ex);
 \fill (1,0) node[above]{1} circle (0.2ex);
 \fill (2,0) node[above]{2} circle (0.2ex);
 \fill (3,0) node[above]{3} circle (0.2ex);
 \fill (4,0) node[above]{4} circle (0.2ex);
 \fill (5,0) node[above]{5} circle (0.2ex);
 \fill (6,0) node[above]{6} circle (0.2ex);
 \fill (7,0) node[above]{7} circle (0.2ex);
 \fill (8,0) node[above]{8} circle (0.2ex);
 \fill (9,0) node[above]{9} circle (0.2ex);
 \fill (0,-1) node[left]{10} circle (0.2ex);
 \fill (1,-1) circle (0.2ex);
 \fill (2,-1) circle (0.2ex);
 \fill (3,-1) circle (0.2ex);
 \fill (4,-1) circle (0.2ex);
 \fill (5,-1) circle (0.2ex);
 \fill (6,-1) circle (0.2ex);
 \fill (7,-1) circle (0.2ex);
 \fill (8,-1) circle (0.2ex);
 \fill (9,-1) circle (0.2ex);
 \fill (0,-2) node[left]{20} circle (0.2ex);
 \fill (1,-2) circle (0.2ex);
 \fill (2,-2) circle (0.2ex);
 \fill (3,-2) circle (0.2ex);
 \fill (4,-2) circle (0.2ex);
 \fill (5,-2) circle (0.2ex);
 \fill (6,-2) circle (0.2ex);
 \fill (7,-2) circle (0.2ex);
 \fill (8,-2) circle (0.2ex);
 \fill (9,-2) circle (0.2ex);
 \fill (0,-3) node[left]{30} circle (0.2ex);
 \fill (1,-3) circle (0.2ex);
 \fill (2,-3) circle (0.2ex);
 \fill (3,-3) circle (0.2ex);
 \fill (4,-3) circle (0.2ex);
 \fill (5,-3) circle (0.2ex);
 \fill (6,-3) circle (0.2ex);
 \fill (7,-3) circle (0.2ex);
 \fill (8,-3) circle (0.2ex);
 \fill (9,-3) circle (0.2ex);
 \fill (0,-4) node[left]{40} circle (0.2ex);
 \fill (1,-4) circle (0.2ex);
 \fill (2,-4) circle (0.2ex);
 \fill (3,-4) circle (0.2ex);
 \fill (4,-4) circle (0.2ex);
 \fill (5,-4) circle (0.2ex);
 \fill (6,-4) circle (0.2ex);
 \fill (7,-4) circle (0.2ex);
 \fill (8,-4) circle (0.2ex);
 \fill (9,-4) node[right]{49} circle (0.2ex);
 \fill (0,-5) node[left]{50} circle (0.2ex);
 \fill (1,-5) circle (0.2ex);
 \fill (2,-5) circle (0.2ex);
 \fill (3,-5) circle (0.2ex);
 \fill (4,-5) circle (0.2ex);
 \fill (5,-5) circle (0.2ex);
 \fill (6,-5) circle (0.2ex);
 \fill (7,-5) circle (0.2ex);
 \fill (8,-5) circle (0.2ex);
 \fill (9,-5) circle (0.2ex);
 \fill (0,-6) node[left]{60} circle (0.2ex);
 \fill (1,-6) circle (0.2ex);
 \fill (2,-6) circle (0.2ex);
 \fill (3,-6) circle (0.2ex);
 \fill (4,-6) circle (0.2ex);
 \fill (5,-6) circle (0.2ex);
 \fill (6,-6) circle (0.2ex);
 \fill (7,-6) circle (0.2ex);
 \fill (8,-6) circle (0.2ex);
 \fill (9,-6) circle (0.2ex);
 \fill (0,-7) node[left]{70} circle (0.2ex);
 \fill (1,-7) circle (0.2ex);
 \fill (2,-7) circle (0.2ex);
 \fill (3,-7) circle (0.2ex);
 \fill (4,-7) circle (0.2ex);
 \fill (5,-7) circle (0.2ex);
 \fill (6,-7) circle (0.2ex);
 \fill (7,-7) circle (0.2ex);
 \fill (8,-7) circle (0.2ex);
 \fill (9,-7) circle (0.2ex);
 \fill (0,-8) node[left]{80} circle (0.2ex);
 \fill (1,-8) circle (0.2ex);
 \fill (2,-8) circle (0.2ex);
 \fill (3,-8) circle (0.2ex);
 \fill (4,-8) circle (0.2ex);
 \fill (5,-8) circle (0.2ex);
 \fill (6,-8) circle (0.2ex);
 \fill (7,-8) circle (0.2ex);
 \fill (8,-8) circle (0.2ex);
 \fill (9,-8) circle (0.2ex);
 \fill (0,-9) node[left]{90} circle (0.2ex);
 \fill (1,-9) circle (0.2ex);
 \fill (2,-9) circle (0.2ex);
 \fill (3,-9) circle (0.2ex);
 \fill (4,-9) circle (0.2ex);
 \fill (5,-9) circle (0.2ex);
 \fill (6,-9) circle (0.2ex);
 \fill (7,-9) circle (0.2ex);
 \fill (8,-9) circle (0.2ex);
 \fill (9,-9) node[right]{99} circle (0.2ex);
% \draw (4,-7)--(5,-7)--(5,-6)--(6,-6)--(5,-7)--(6,-7)--(8,-9)--(1,-9)--(4,-7)--(5,-8)--(7,-6)--(7,-5)--(6,-3)--(4,-2)--(3,-3)--(4,-6)--(4,-7)--(3,-7)--(1,-6)--(1,-4)--(2,-2)--(3,-1)--(4,-1)--(6,0)--(6,-1)--(4,0)--(4,-1)--(7,-1)--(8,-2)--(7,-2)--(7,-3)--(8,-5)--(8,-6)--(7,-6);
\end{tikzpicture}\end{center}

\end{document}
